\section{Analysis}
\label{sec:analysis}

\subsection{Error Taxonomy}
\label{sec:error-taxonomy}

Under the current graph construction, all observed end-to-end failures (recall = 0) in graph-based routing originate from incorrect type prediction rather than graph coverage gaps.
The oracle strategy achieves F1 = 1.000 across all five domains, confirming that every query has a valid BFS path when ground-truth types are provided.

This means the graph is never the bottleneck: when the LLM predicts the correct entity types, the system always produces the correct tool sequence.
Improvement efforts should therefore focus on type prediction quality rather than graph coverage.

\subsection{Regressions}
\label{sec:regressions}

We observe 50 query instances across all model--domain combinations where the baseline achieves non-zero recall but graph routing achieves zero.
All 50 are caused by type prediction failures: the LLM predicts incorrect source or target types, so BFS either finds no path or finds a path to the wrong tools.

This tradeoff is intentional.
The graph will not hallucinate tools---it returns nothing rather than something wrong.
Whether this is desirable depends on the application.
In high-stakes domains (infrastructure operations, financial systems), returning nothing is safer than invoking wrong tools.
In exploratory domains, partial results from the baseline may be more useful.

\subsection{Model Effect}
\label{sec:model-effect}

In our experiments, smaller models exhibited larger improvements from graph routing.
Granite 8B averages $\Delta$F1 = +0.406 across domains, while Qwen 14B averages +0.264.
This pattern is consistent with the hypothesis that graph constraints compensate for weaker type prediction by reducing the decision space, but we note this observation is based on four models and should not be interpreted as a general law.

\subsection{Mechanistic Explanation}
\label{sec:mechanistic}

We provide a quantitative account of \emph{why} graph routing improves performance by measuring how the graph reduces the decision space at each stage.

\paragraph{Tool Space Reduction.}
Graph-based routing reduces the tool candidate set by 97--99\% compared to the full catalog.
The BFS path presents an average of 1.4--2.0 tools per query, compared to 54--170 tools in the full catalog (Table~\ref{tab:domain-stats}).
This is an extreme form of context filtering: the model sees only the tools that are structurally valid for the predicted types.

\paragraph{Entity Space Reduction.}
Entity pruning (Section~\ref{sec:entity-pruning}) removes 56--87\% of entity types as source candidates after target prediction.
In CI/CD, where 51 entity types exist, a typical query constrains source prediction to at most 7 candidates (87\% pruning).
This is a guaranteed bound, independent of model quality.

\paragraph{Structural Validity.}
Every tool sequence returned by graph routing is compositionally valid by construction: the output type of each tool matches the input type of the next.
The baseline has no such guarantee, which accounts for its hallucinated tool invocations (Table~\ref{tab:hallucinations}).

Together, these three mechanisms explain the performance improvement: the graph reduces the decision space from $O(|\mathcal{T}|)$ to $O(|\text{path}|)$, eliminates structurally invalid options, and constrains each prediction step through reachability.
The LLM makes a simpler decision (choose among $|\mathcal{E}|$ types, further constrained by reachability) rather than a harder one (choose among $|\mathcal{T}|$ tools with no structural guidance).

\subsection{Fault Tolerance}
\label{sec:fault-tolerance}

The recall decomposition (Section~\ref{sec:recall-decomposition}) reveals that the graph is resilient to type prediction errors.
With $R_{\text{wrong}}$ averaging 0.415, incorrect predictions recover roughly 40\% of the correct tool chain.
This occurs when the predicted type is a graph neighbor of the correct type: BFS from a nearby node traverses through some of the same edges.

This property is desirable: it means that minor type prediction errors (predicting a parent or sibling type) degrade gracefully rather than producing total failure.
The graph's topological structure provides a soft landing for prediction mistakes.
